\section{What The Chapter Carries Forward}

The chapter began with the invariant left by the finite gauge. It has now
grounded that invariant in a concrete operator. A single clock reads one
burden in two ways: answerability and elapsed cost. A boundary-anchored
chain gives the finite Poincare reading: zero energy forces zero activity.
The square-root door turns energy squared into a scale field fit for
completion. The universe kernel binds labels to the invariant without
letting labels become independent measures. Metaphysical relativity states
that an invariant is metaphysical relative to a sensor class that cannot
resolve it and physical relative to a class that can carry it. Finally, the
operator instantiates the derivative pipeline itself.

What passes forward is therefore not merely a useful example. It is a
grounded grammar of activity. The next chapter may read signs, pairs, and
geometric turns, but it inherits the discipline established here: activity
is measured by one positive invariant, zero activity is detected by the
anchored operator, labels ride the invariant, and physical record depends on
the resolving power of the sensor class.
